%% kapa: exam=1 round=21 year=2010 period=1 form=A subject=civil_law subject_ko=민법 seq=12 number=12
%% source: Q-Net 감정평가사 기출문제 (cst00309) artlSeq=5024181 file=extracted/1008/5024181/2010년 제21회 감정평가사 1차시험 기출문제/1교시.hwp
%% conversion: kordoc 4.24.51 -> tools/gen1.py (auto); review=proofread-sonnet; flags=none
\begin{question}{12}
부동산 등기에 관한 설명으로 옳지 않은 것은? (다툼이 있으면 판례에 의함)
\begin{choices}
  \item 부기등기의 순위는 주등기의 순위에 의한다.
  \item 경정등기는 등기와 실체관계 사이의 후발적 불일치를 시정하기 위한 등기이다.
  \item 동일한 부동산에 관하여 등기한 권리의 순위는 법률에 다른 규정이 없으면 등기한 순서에 따른다.
  \item 지상권의 존속기간에 관한 등기는 제3자에 대하여 대항력이 있다.
  \item 가등기에 기초한 본등기가 경료되기 전에는 가등기된 청구권의 기초인 법률관계가 존재한다는 추정력이 인정되지 아니한다.
\end{choices}
\end{question}
\answer{2}
